\documentclass{../../slides-style}

\usepackage{booktabs}

\slidetitle{Профилирование производительности}{28.02.2026}

% (c) Николай Пономарев

\subject{Использование инструментов сбора и анализа данных о производительности на примере perf tools.  }

\begin{document}

\begin{frame}[plain, noframenumbering]
    \titlepage
\end{frame}

\section{Мотивация}

\begin{frame}{Зачем это вообще нужно}
    \begin{itemize}
        \item Если вдруг программа работает медленно...
        \item Можно угадать, где проблема
        \item Можно измерить, где проблема~--- заняться \textbf{профилированием}
        \item Иногда хватит первого, но второе работает чаще
    \end{itemize}
\end{frame}

\begin{frame}{Пример этой пары}
    \begin{columns}
        \column{0.48\textwidth}
        \begin{block}{Программа A}
            \texttt{sum += a[\alert{i}][\alert{j}];}
        \end{block}
        \column{0.48\textwidth}
        \begin{block}{Программа B}
            \texttt{sum += a[\alert{j}][\alert{i}];}
        \end{block}
    \end{columns}

    \bigskip
    \begin{itemize}
        \item \texttt{double a[4096][4096]}
        \item Одинаковое число операций: $4096^2$ сложений
        \item Одинаковый результат
        \item Спойлер: разница в скорости --- в несколько раз
        \item Будем разбираться
    \end{itemize}
\end{frame}

\section{Иерархия памяти}

\begin{frame}{Cache line}
    \begin{itemize}
        \item Процессор не подгружает один байт
        \item Единица обмена --- \textbf{64 байта} (cache line, линия кеша)
        \item При доступе к \enquote{близкому} элементу~--- обратимся в кеш
        \item При доступе к \enquote{далекому} элементу~--- пойдём в ОЗУ, а это дорого
    \end{itemize}

    \bigskip
    Spatial locality --- работает быстрее, потому что кэш

    Cache miss --- работает медленнее, потому что RAM

\end{frame}

\begin{frame}{Иерархия памяти}
    \begin{center}
        \begin{scriptsize}
            \begin{tabular}{lrrr}
                \toprule
                Уровень  & Размер (типично) & Латентность (такты) & Латентность (нс) \\
                \midrule
                Регистры & $<$ 1 КБ         & $\sim$1             & $<$1             \\
                L1 cache & 32--64 КБ        & $\sim$4--5          & $\sim$1--2       \\
                L2 cache & 256 КБ -- 2 МБ   & $\sim$10--15        & $\sim$3--5       \\
                L3 cache & 4--64 МБ         & $\sim$30--50        & $\sim$10--20     \\
                RAM      & ГБ               & $\sim$150--300      & $\sim$50--100    \\
                \bottomrule
            \end{tabular}
        \end{scriptsize}
    \end{center}

    \smallskip
    {\tiny Типичные значения для современных x86 процессоров; конкретные цифры зависят от микроархитектуры и частоты}

    \bigskip

    Пока процессор ждёт данные из RAM, он мог бы выполнить $\sim$200 инструкций
\end{frame}

\begin{frame}{Как лежит массив в памяти}
    \texttt{double a[M][N]} в C хранится строка за строкой (row-major):

    \bigskip
    \begin{center}
        \ttfamily\small
        {\renewcommand{\arraystretch}{1.2}
            \begin{tabular}{|c|c|c|c|c|c|c|c|}
                \hline
                a[0][0]                                  & a[0][1]                                 & a[0][2] & a[0][3] & a[1][0] & a[1][1] & a[1][2] & \ldots \\
                \hline
                \multicolumn{4}{|c|}{\rmfamily строка 0} & \multicolumn{4}{c|}{\rmfamily строка 1}                                                            \\
                \hline
            \end{tabular}
        }
    \end{center}

    \bigskip
    \begin{columns}
        \column{0.48\textwidth}
        \begin{block}{\texttt{a[i][j]}, внешний \texttt{i}}
            Последовательный доступ\\
            Cache-friendly
        \end{block}
        \column{0.48\textwidth}
        \begin{alertblock}{\texttt{a[j][i]}, внешний \texttt{j}}
            Прыжок на $N \times 8$ байт\\
            При $N=4096$: \textbf{32~768 байт}\\
            Cache-hostile
        \end{alertblock}
    \end{columns}
\end{frame}

\section{perf}

\begin{frame}{perf}
    \begin{itemize}
        \item Стандартный инструмент профилирования в Linux
        \item Использует PMU (Performance Monitoring Unit) --- счётчики прямо в железе
        \item Низкий overhead, не замедляет программу принципиально
    \end{itemize}

    \bigskip
    \begin{center}
        \begin{tabular}{ll}
            \toprule
            \texttt{perf stat}     & сводные счётчики по всей программе        \\
            \texttt{perf record}   & записать профиль с сэмплингом             \\
            \texttt{perf report}   & интерактивный TUI с hot-spot'ами          \\
            \texttt{perf annotate} & ассемблерный листинг с \% времени         \\
            \texttt{perf top}      & аналог top для анализа производительности \\
            \bottomrule
        \end{tabular}
    \end{center}
\end{frame}

\begin{frame}{Что именно считает perf}
    \begin{itemize}
        \item \textbf{Hardware events} --- считаются прямо в железе (PMU)
              \begin{itemize}
                  \item \texttt{cycles} --- тактов процессора
                  \item \texttt{instructions} --- выполненных инструкций
                  \item \texttt{cache-misses} --- промахов мимо кэша
                  \item \texttt{branch-misses} --- неверно предсказанных ветвлений
              \end{itemize}

              \bigskip
        \item \textbf{Software events} --- считаются ядром
              \begin{itemize}
                  \item \texttt{task-clock} --- реальное время CPU
                  \item \texttt{page-faults} --- промахов в таблице страниц
              \end{itemize}

              \bigskip
        \item \textbf{Tracepoints} --- системные вызовы, планировщик, сеть
    \end{itemize}

    \bigskip

    \textbf{IPC} (Instructions Per Cycle) --- сколько инструкций за такт

    IPC $< 1$ --- процессор ждёт данные
\end{frame}

\section{Демонстрация}

\begin{frame}[fragile]{Программы для анализа}
    \begin{columns}[T]
        \column{0.48\textwidth}
        \inputminted[]{c}{row.c}
        \column{0.48\textwidth}
        \inputminted[]{c}{col.c}
    \end{columns}
\end{frame}

\begin{frame}[fragile]{Собираем и запускаем}
    \begin{minted}{console}
$ gcc -O0 -g -o row row.c
$ gcc -O0 -g -o col col.c

$ /usr/bin/time ./row
$ /usr/bin/time ./col
\end{minted}

    \begin{itemize}
        \item Просто \texttt{time} может быть встроенной функцией оболочки
        \item Строго говоря, измерения~--- отдельная наука, об этом расскажут
        \item Начать упрощать себе жизнь можно с hyperfine: \url{https://github.com/sharkdp/hyperfine}
    \end{itemize}

\end{frame}

\begin{frame}[fragile]{perf stat}
    \begin{minted}{console}
$ perf stat ./row
$ perf stat ./col
\end{minted}

    \bigskip
    \begin{center}
        \begin{tabular}{lll}
            \toprule
            Метрика               & row          & col                      \\
            \midrule
            \texttt{task-clock}   & $\sim$500 мс & $\sim$2500 мс            \\
            \texttt{instructions} & $\approx$ X  & $\approx$ X              \\
            \texttt{cache-misses} & мало         & \alert{\textbf{огромно}} \\
            \bottomrule
        \end{tabular}
    \end{center}

\end{frame}

\begin{frame}[fragile]{perf record + perf report}
    \begin{minted}{console}
$ perf record -g ./col
$ perf report
\end{minted}

    \bigskip
    В TUI:
    \begin{itemize}
        \item \textbf{Enter} на \texttt{main} --- раскрыть call graph
        \item \textbf{a} --- ассемблерный листинг
        \item \texttt{movsd} --- загрузка \texttt{double} из памяти --- самое \enquote{горячее} место
    \end{itemize}

\end{frame}

\begin{frame}[fragile]{Flame Graphs}
    \begin{itemize}
        \item Некоторое развитие \texttt{perf report}
        \item Визуализация профиля: ось X --- время, ось Y --- стек вызовов
        \item Широкий прямоугольник --- много времени в этой функции
        \item Позволяет быстро найти горячие пути в больших программах
    \end{itemize}

    \bigskip
    \begin{minted}{console}
$ perf record -g ./col
$ perf script | stackcollapse-perf.pl | flamegraph.pl > out.svg
\end{minted}

    \bigskip
    \begin{itemize}
        \item Скрипты: \url{https://github.com/brendangregg/FlameGraph}
        \item Современный вариант без Perl: \texttt{inferno} (Rust)
    \end{itemize}
\end{frame}

\section{Выводы}

\begin{frame}[fragile]{А если собрать с -O2?}
    \begin{minted}{console}
$ gcc -O2 -g -o col_opt col.c
$ perf stat ./col_opt
$ perf stat ./row       # -O0!
\end{minted}

    \bigskip
    \begin{itemize}
        \item \texttt{col\_opt} быстрее \texttt{col}
        \item Но медленнее \texttt{row} без всяких оптимизаций
    \end{itemize}

    \bigskip

    Компилятор не знает паттерн доступа к памяти, а вот программист...

\end{frame}

\begin{frame}{Типичные ошибки при профилировании}
    \begin{itemize}
        \item Профилируйте c хотя бы \texttt{-O2 -g}
              \begin{itemize}
                  \item В CMake используйте \texttt{RelWithDebInfo}
              \end{itemize}

              \bigskip
        \item Запускайте много раз

              \bigskip
        \item Корректность алгоритма важнее скорости

    \end{itemize}
\end{frame}

\begin{frame}{Что оптимизировать и в каком порядке}
    \begin{enumerate}
        \item \textbf{Компилятор} (\texttt{-O2}, \texttt{-O3}, \texttt{-march=<...>})
              \begin{itemize}
                  \item Бесплатно, включайте всегда
                  \item Loop unrolling, инлайнинг, векторизация
              \end{itemize}

              \bigskip
        \item \textbf{Алгоритм и структуры данных}
              \begin{itemize}
                  \item Порядки величин
                  \item Cache-friendly vs cache-hostile
              \end{itemize}

              \bigskip
        \item \textbf{Микрооптимизации} (SIMD, intrinsics)
              \begin{itemize}
                  \item Проценты
                  \item Сложность реализации кратно выше
              \end{itemize}
    \end{enumerate}
\end{frame}

\begin{frame}{Полезные ссылки}
    \begin{itemize}
        \item Большой практический гид по perf с примерами
              \begin{itemize}
                  \item \url{https://brendangregg.com/perf.html}
              \end{itemize}

        \item \texttt{man perf-stat}, \texttt{man perf-record}
    \end{itemize}
\end{frame}


\end{document}
